% ============================================================
% 01_PROBLEMA.TEX
% Sistema Adaptativo de Classificação Multicritério
% de Fluidos de Corte
%
% FINALIDADE
%
% Apresentar o problema de decisão que motiva o projeto.
%
% Este arquivo constitui apenas a BASE da apresentação.
% Os integrantes deverão posteriormente complementar a
% contextualização com as referências e informações técnicas
% efetivamente utilizadas no trabalho.
%
% Não inserir resultados numéricos nesta seção.
% ============================================================


% ============================================================
% FRAME 1 — PROBLEMA DE DECISÃO
% ============================================================

\begin{frame}{Problema de decisão}

    \begin{block}{Pergunta central}
        Como classificar os fluidos de corte de forma
        \textbf{transparente, reprodutível e tecnicamente justificável},
        considerando simultaneamente múltiplos critérios de avaliação?
    \end{block}

    \vspace{0.3cm}

    A comparação entre alternativas não deve depender de uma única
    variável, pois diferentes fluidos podem apresentar vantagens e
    desvantagens em diferentes dimensões.

    \vspace{0.3cm}

    \begin{itemize}
        \item uma alternativa pode apresentar bom desempenho em um critério
              e desempenho inferior em outro;

        \item critérios podem possuir unidades, escalas e direções distintas;

        \item algumas informações podem representar desempenho, enquanto
              outras podem representar restrições ou conformidade técnica;

        \item inconsistências ou problemas na base de dados podem afetar
              diretamente qualquer classificação posterior.
    \end{itemize}

\end{frame}


% ============================================================
% FRAME 2 — POR QUE MULTICRITÉRIO?
% ============================================================

\begin{frame}{Por que uma abordagem multicritério?}

    \begin{columns}[T]

        % ----------------------------------------------------
        % COLUNA ESQUERDA
        % ----------------------------------------------------

        \begin{column}{0.48\textwidth}

            \begin{block}{Problema}
                Avaliar cada variável isoladamente pode produzir
                conclusões diferentes para a mesma alternativa.
            \end{block}

            \vspace{0.2cm}

            Portanto, é necessário considerar conjuntamente:

            \begin{itemize}
                \item desempenho;
                \item qualidade;
                \item aspectos técnicos;
                \item aspectos econômicos, quando disponíveis;
                \item estabilidade entre condições;
                \item conformidade.
            \end{itemize}

        \end{column}


        % ----------------------------------------------------
        % COLUNA DIREITA
        % ----------------------------------------------------

        \begin{column}{0.48\textwidth}

            \begin{block}{Necessidade}
                A classificação deve combinar os critérios sem ocultar
                os compromissos existentes entre eles.
            \end{block}

            \vspace{0.2cm}

            O projeto deverá permitir:

            \begin{itemize}
                \item identificar redundâncias;
                \item ponderar critérios de forma reproduzível;
                \item comparar as alternativas;
                \item verificar dominância;
                \item avaliar sensibilidade;
                \item avaliar robustez do ranking.
            \end{itemize}

        \end{column}

    \end{columns}

\end{frame}


% ============================================================
% FRAME 3 — PRINCÍPIO DO PROJETO
% ============================================================

\begin{frame}{Princípio do projeto}

    \begin{center}

        \textbf{Não classificar antes de compreender e validar os dados.}

    \end{center}

    \vspace{0.4cm}

    \begin{block}{Sequência lógica}
        \centering
        Dados
        $\rightarrow$
        Auditoria
        $\rightarrow$
        Critérios
        $\rightarrow$
        Ponderação
        $\rightarrow$
        Classificação
        $\rightarrow$
        Robustez
    \end{block}

    \vspace{0.4cm}

    Assim, o ranking final deverá ser consequência de uma cadeia
    metodológica rastreável, e não de uma escolha direta baseada
    apenas em observações isoladas da planilha.

    \vspace{0.3cm}

    \begin{alertblock}{Importante}
        Nesta fase de preparação da estrutura, nenhum fluido é definido
        como superior aos demais. Essa conclusão deverá surgir somente
        após a execução das análises pelos integrantes.
    \end{alertblock}

\end{frame}


% ============================================================
% MATERIAL EXTERNO
% ============================================================
%
% Caso a professora forneça posteriormente:
%
% - definição específica do problema;
% - texto de contextualização;
% - requisito técnico;
% - referência;
% - imagem;
%
% o conteúdo deverá ser incorporado preservando a fonte e
% sem substituir silenciosamente o material oficial.
%
% ============================================================


% ============================================================
% ESTADO ATUAL
% ============================================================
%
% PROBLEM_SECTION_STATUS = BASE_PREPARED
%
% Nenhum resultado científico foi produzido nesta etapa.
%
% ============================================================